% Declaración de uso de IA — Responsable: Apunte Pazmiño Washington German
% Tiene que decir lo que realmente pasó. Si alguien aclara otro uso de IA, se agrega aquí.
\section*{Declaración de uso de inteligencia artificial}

\begingroup\footnotesize
\textbf{Herramientas de apoyo:} Claude (Anthropic; modelos Opus 5.5 y Sonnet 5.5), consultado el 14 y 15 de noviembre de 2026, y ChatGPT (OpenAI), utilizado como asistente por Jordan Loor.

\textbf{Uso como herramienta de asistencia en el trabajo del equipo:}
\begin{itemize}
\item \textbf{Apoyo en verificación bibliográfica:} asistencia para contrastar metadatos en Crossref (autores, título, revista, páginas y DOI) a partir de las fuentes localizadas por el equipo.
\item \textbf{Revisión de estilo y maquetación:} sugerencias de corrección ortográfica, gramatical y de sintaxis en formato LaTeX.
\item \textbf{Estructuración técnica:} apoyo en la organización de la plantilla LaTeX, el archivo README y la estandarización de comandos de git para el flujo de trabajo del equipo.
\item \textbf{Asistencia de formato:} apoyo en la depuración de etiquetas XML para el diagrama de Ishikawa en draw.io y revisión preliminar de estructura para los apartados 2 y 3.
\end{itemize}

\textbf{Autoría y participación de cada integrante:}
\begin{itemize}
\item \textbf{Washington Apunte:} realizó la lectura crítica de las fuentes y redactó íntegramente el apartado 4. Utilizó Claude únicamente como apoyo para revisar el estilo de redacción, contrastar cifras contra el informe oficial y depurar el formato LaTeX de las citas.
\item \textbf{Jordan Loor:} investigó y redactó los apartados 1, 5 y 6, así como las observaciones. Empleó ChatGPT como herramienta de consulta para pulir la cohesión de la descripción del incidente, validar la estructura de las medidas de mitigación y riesgos de SREI definidos por él, y verificar la sintaxis LaTeX antes de integrar su aporte.
\item \textbf{César Macías:} desarrolló el análisis y el contenido de los apartados 2 y 3, así como el diseño conceptual del diagrama de Ishikawa, apoyándose en la IA para optimizar la sintaxis del código fuente (XML) en draw.io y la estructura base. Ejecutó manualmente la exportación de los archivos \texttt{.drawio} y \texttt{.pdf}, la limpieza y organización del repositorio y las correcciones técnicas en LaTeX (SLA Business y la precisión sobre el ciberataque).
\end{itemize}
\endgroup